\subsection{Model validation and robustness testing through targeted perturbations}

To validate our computational model and assess network robustness, we reproduced key perturbation experiments from the original \textit{Dendronotus iris} swim CPG studies \cite{SakuraiKatz2016, sakurai2022bursting}.
Brief hyperpolarizing or depolarizing current pulses delivered to targeted interneurons during ongoing swim episodes test how the circuit responds to transient disruptions and whether it recovers its characteristic oscillatory pattern.
A robust network should maintain coordinated output despite perturbations and return rapidly to stable operation upon their removal.
Faithful reproduction of these experimental responses confirms that the model captures both the functional dynamics and the stability properties of the biological CPG.

Figure~\ref{fig:perturbation_responses} presents six perturbation experiments (panels A--F) that reveal distinct aspects of circuit resilience.
Across all conditions, the circuit recovers coordinated bursting once perturbations are removed, demonstrating that the network operates as a robust attractor.

Panel A shows that depolarizing Si3L produces immediate cessation of network bursting followed by rapid recovery.
The contralateral Si3R neuron becomes hyperpolarized due to strong reciprocal inhibition, and Si2R is silenced via the E/I inhibitory projection.
Panel B represents a special case: it is the only perturbation in which network bursting is not completely halted.
The targeted Si2R neuron is strongly inhibited, yet activity persists through alternative pathways---underscoring that the E/I modules, not individual neurons, constitute the core oscillatory engine.

Panel C shows both Si2 interneurons simultaneously hyperpolarized.
The circuit recovers, but restoration takes noticeably longer than in unilateral cases, highlighting the supporting role Si2 neurons play in stabilizing the rhythm.
Panel D applies a depolarizing pulse to Si3R, disrupting alternation; upon pulse termination, the circuit rapidly returns to its characteristic antiphase rhythm.

Panel E shows hyperpolarization of Si3R, which silences its electrical partner Si2L through gap junction coupling.
Panel F shows bilateral hyperpolarization of both Si3 neurons, which completely abolishes network activity.
Silencing both Si3 neurons eliminates excitatory drive to the Si2 neurons, which are configured near the threshold between tonic spiking and quiescence.
Upon removal of the hyperpolarizing pulses, the CPG rapidly restores the control rhythm, demonstrating that Si3 neurons are necessary to initiate network-level bursting.
Additional symmetric bilateral perturbation experiments further validate the model's dynamics (Supplementary Materials).

These perturbation experiments assumed elevated Si3 excitability to sustain ongoing bursting.
In the biological circuit, this drive originates from Si1 command neurons, whose control over network state we examine next.
